\scanpage{35}\mspage{34}
% L1–2
Soit $\deleted{X}A\in M$. On va ``relativiser'' \unsure{tout} p.r. à $\deleted{X}A$. Soit \deleted{$M_X=M_{\leq X}$,}
% L3
\[
M_{\deleted{X}A}=\deleted{\unsure{\{}}M_{\leq\deleted{X}A}=\{Y\in M\mid Y\subset\deleted{X}A\}
\]
% L4–5
$M_{\deleted{X}A}$ muni de la relation d'ordre induite par $M$

% L6
$\{X,\deleted{\unsure{Y}}Z\}$ est mod dans $M_{\deleted{X}A}$ \unsure{sss|ssi} il l'est dans $M$

% L7–9
\[
\begin{aligned}
X\in M_{\deleted{\unsure{X}}A}\text{ est fermé}\quad &\text{sss}\qquad X=\overline X\wedge\deleted{X}A\\
&\text{i.e. sss}\quad\exists\,\deleted{\unsure{Y}}Z\in M_f,\ \{\deleted{\unsure{Y}}Z,\deleted{X}A\}\ \text{mod et}\\
&\hspace{10em}X=\deleted{\unsure{Y}}Z\wedge\deleted{X}A
\end{aligned}
\]

% L10–11
Je dis que les conditions 1) à 7) sont encore satisfaites~:
% L12, 1)
\[
\text{1)}\quad\varnothing_{M_{\deleted{\unsure{X}}A}}=\varnothing_{M_{\deleted{\unsure{\leq X|X|A}}}},\qquad 1_{M_{\deleted{\unsure{X}}A}}=X
\]
% L13, 2)
2) \deleted{$\unsure{X\wedge Z}$} trivial

% L14–15, 3)
3) item car \deleted{$\unsure{X\overset{M}{\wedge}Z}$} $X\overset{M}{\wedge}Y$ et $X\overset{M}{\vee}Y$ sont dans $M_A$ donc ce sont des bornes inf et sup dans $M_A$.

% L16–23, 4)
4) \deleted{$X\unsure{<|\subset}A$,} $X\in M_A$, considérons les $Y\in M_{\unsure{Y|y}}$ en position mod. p.r. à $X$, telles que $Y\wedge X=\varnothing$ i.e. \deleted{les $Y\subset\unsure{X'|X}$} on a $X\subset A$, on considère les $Y\subset A$ tels que $\{X,Y\}$ mod. et $X\wedge Y=\varnothing$. \unsure{Ces|Les} $Y$ sont $\leq X'$ (déf de $X'$) et $\leq A$, donc $\leq X'\wedge A$, qui a les propriétés voulues, et c'est \unsure{la|le} plus \unsure{grande|grand} de tous, OK

